\documentclass[11pt]{article}
\usepackage{amsmath,amssymb,amsthm,mathrsfs}
\usepackage[margin=1in]{geometry}
\usepackage[hidelinks]{hyperref}
\newcommand{\DD}{d'd''}
\newcommand{\MA}{\operatorname{MA}}
\newcommand{\C}{\mathcal C}
\newcommand{\G}{\mathcal G}
\newcommand{\RR}{\mathbb R}
\newcommand{\HH}{\mathscr H}
\newtheorem{theorem}{Theorem}
\newtheorem{lemma}[theorem]{Lemma}
\newtheorem{corollary}[theorem]{Corollary}
\newtheorem{proposition}[theorem]{Proposition}
\theoremstyle{remark}\newtheorem{remark}[theorem]{Remark}
\title{Uniqueness for actual decreasing-generated finite-energy potentials\\
Two actual Gauss fibres, point atoms, and bounded test energies}
\author{Research working paper 403; unreviewed}
\date{9 October 2026}
\begin{document}
\maketitle
\begin{abstract}
For the actual nonalgebraic two-dimensional analytic torus quotient,
we give a pointwise uniqueness proof in the genuinely
decreasing-generated finite-energy class. Equal Monge--Ampere
measures force every bounded test gradient energy of the continuous
presentation differences to tend to zero. Positive Gauss-fibre
currents then control differences at curve endpoints. A full
graded-dimensional point is connected to the reference Gauss
section by two actual coordinate fibres, with finite extensions
used only for curve norms. Positive bounded point atoms make the
endpoint limits finite. Reconstruction from all affinoid Shilov
boundaries proves equality up to one constant, including at poles.
No unbounded divisor restriction or surface Monge--Ampere
base-change formula is used. The intrinsic maximal finite-energy
domain, energy-sublevel representability and compactness, and
arbitrary-measure existence remain unresolved. All new claims in
this working paper are unreviewed; original P16.7 is not declared
solved.
\end{abstract}

\section{Object, generated class, and named foundations}
Let $K=\mathbb C((s))$, $|s|=e^{-1/8}$, $t=s^8$, and
\[
Y=(\mathbb G_m^2)^{\rm an}_K,\qquad
X=Y/\langle(t,2),(3,t)\rangle,\qquad
\tau=(-\log|z_1|,-\log|z_2|).
\]
We use the reconstructed compact, separated, good, smooth,
boundaryless analytic quotient. Tropical boxes of width less than
one give ordinary injective quotient charts. Its quotient and
nonalgebraization arguments are standing geometric inputs, not
conclusions of uniqueness. Put $\rho=|\tau|^2/2$ locally and
$\theta=\DD\rho$. Reference changes are smooth affine unit
logarithms. The actual saturated degree-one calibration gives raw
mass $\int_X\theta^2=M=2$.

Let $\mathscr P$ be the minimal continuous CLD psh sheaf closed under
local uniform limits and containing the ordinary smooth psh
functions. All continuous stages are included. Let
\[
\C=\{a\in C(X,\RR):\rho+a\in\mathscr P\text{ locally}\},\quad
\theta_a=\theta+\DD a,\quad B(a,b)=\theta_a\wedge\theta_b,
\]
and
\[
E(a)=\frac13\int a[\theta^2+B(a,0)+B(a,a)].
\]
The actual generated class $\G$ consists of pointwise limits
$a_j\downarrow u$ with $a_j\in\C$ and $\inf_jE(a_j)>-\infty$.
The limits and point values are actual USC functions, not classes
modulo null sets.

The appendix reproduces the whole point-atom and bounded Green
argument of working paper 402. In particular, at every continuous
stage there is a positive Hilbert pairing
\[
Q_T(f,g)=-\int f\DD(gT)
\]
for positive closed $T$ and continuous dpsh $f,g$, realized by
uniform smooth approximations with strongly converging gradients.
Bounded Green gives the cocycle, first variation and concavity
\begin{align}
E(a)-E(b)&=\frac13\int(a-b)[\MA(a)+B(a,b)+\MA(b)],\label{403cocycle}\\
\frac d{dt}E((1-t)a+tb)&=\int(b-a)\MA((1-t)a+tb),\label{403var}\\
\frac{d^2}{dt^2}E((1-t)a+tb)&=-2Q_{\theta_{(1-t)a+tb}}(b-a,b-a)\le0.
\end{align}
These use actual recursive CLD continuous products, positivity,
locality and local uniform continuity. The named source inputs are
CLDv2 7.2.8, 8.5.1, 8.5.3, 8.5.5--6, 8.7.1--4, 8.8.4--5,
and ordinary smooth Stokes, cutoffs and scalar approximation.
Arbitrary continuous-multiplier closedness in 8.6.7 is not used.

For Gauss-fibre currents we use CLD 18.6.2 with its 18.6.1
relative curvature formula, at the boundaryless, separated,
equidimensional hypotheses. Their immediate proofs were read;
their relative intersection foundations in 18.3--18.4 are named
published inputs, not new recursively proved results. We use
14.5.2--3 for canonical curve calibration, 15.3.1 and 16.1.1 for
normalized varying-ground-field integration, and 14.4.3 for the
different operation of viewing a finite-extension curve over the
old ground field. The full applied immediate units were read.
Standard analytic scalar extension, smooth pullback, the
algebraic closedness of the completed algebraic closure and the
graded Abhyankar inequality are standing analytic foundations.
The disc classification needed below is proved explicitly.

\section{The canonical generated operator and all four moving moments}
We include these steps because convergence against fixed continuous
tests alone does not justify a moving unbounded moment.
Compact strict domination makes $E(u)=\lim_j E(a_j)$ independent of
every continuous decreasing presentation: if $b_k\downarrow u$
and continuous $a>u$, the increasing ordinary open sets $\{b_k<a\}$
cover compact $X$, so one index works globally. Compare with
$a_j+\epsilon$, apply order and $E(a+c)=E(a)+Mc$, and reverse
the presentations. The same argument proves generated order
preservation and continuity when a decreasing generated sequence
already has a generated limit. Convex combinations are generated
by the corresponding continuous diagonal presentations.

On a finite simplex of generated weights, bounded energies are
polynomials of degree at most three. Their values converge at a
finite unisolvent set by the preceding energy argument. The
coefficients therefore converge and give a canonical energy
polynomial on the whole simplex. Such an interpolation set exists
because a nonzero polynomial cannot vanish on an open simplex;
zero-weight slots are omitted at boundary points.

For a presentation $a_j\downarrow u$ and a bounded $v\in\C$,
set $P_j(t)=E((1-t)a_j+tv)$ and $R_j(t)=E(ta_j)$.
Equation~\eqref{403var} gives
\[
\int v\MA(a_j)=P'_j(0)+R'_j(1).
\]
The coefficient limits are presentation-independent. The span of
$\C$ is a uniformly dense vector lattice in $C(X)$: tropical
periodic perturbations separate different tropical points, while
small globally capped Laurent polynomial tests separate distinct
seminorms in one tropical fibre. The cap is
$\max(\rho_a,\delta+\epsilon(\log|P|-b))-\rho_a$ on a small
injective box, extended by zero; the zero locus is in its zero
branch and a fixed boundary collar is zero. Finite maxima in
$\C$ give the vector lattice identity
\[
\max(A(p-q),A(r-s))=
2A\max((p+s)/2,(r+q)/2)-A(q+s).
\]
Constants and point separation then give density by the usual
finite-cover lattice approximation. Positivity and the common
mass $M$ extend convergence to every continuous test. Compact
Hausdorff Riesz representation gives the actual positive Radon
limit $\MA(u)$ of mass $M$. Polarization gives the actual positive
mixed limits. This is a limit of bounded analytic products, not
an inverse constructed from a dual functional.

Define $E^*(\mu)=\sup_{a\in\C}(E(a)-\int a\,d\mu)$.
Bounded tangent gives $E^*(\MA(a))=E(a)-\int a\MA(a)$.
The coefficient $D_u=\lim_j\int a_j\MA(a_j)=R'(1)$ exists.
Passing the bounded tangent inequality with a fixed bounded
competitor gives
\[
E^*(\MA(u))\le E(u)-D_u<\infty.
\]
This establishes integrability before any unbounded moving
integral is invoked: for $u\le0$, monotone convergence from a
common continuous upper bound gives
\[
\int(-u)\,d\mu\le E^*(\mu)-E(u)
\quad\text{if }E^*(\mu)<\infty,
\]
and every presentation converges in $L^1(\mu)$.

For $k\ge j$, equation~\eqref{403cocycle} and positivity imply
\[
0\le\int(a_j-a_k)\MA(a_k)\le3(E(a_j)-E(a_k)).
\]
First let $k\to\infty$ using the fixed continuous test $a_j$,
then $j\to\infty$ using the established $L^1(\MA(u))$
convergence. It follows that $D_u=\int u\MA(u)$ and
\[
E^*(\MA(u))=E(u)-\int u\MA(u).
\]

For generated $u,v$ the coefficient of $s$ at $(1,0)$ in
$F(t,s)=E(tu+sv)$ is exactly $\int v\MA(u)$.
For bounded $v$ this is the bounded-test coefficient identity.
For general $v$, decrease bounded $v_l$ to $v$. At every simplex
point the limit $tu+sv$ is already generated, so energy
continuity and polynomial interpolation give coefficient
convergence; finite full dual gives the same limit by monotone
convergence of $\int v_l\MA(u)$. Thus for any two continuous
decreasing presentations $a_j\downarrow u$, $b_j\downarrow v$,
coefficient convergence of $E(ta_j+sb_j)$ proves all four limits
\begin{align*}
\int a_j\MA(a_j)&\longrightarrow\int u\MA(u),&
\int b_j\MA(a_j)&\longrightarrow\int v\MA(u),\\
\int a_j\MA(b_j)&\longrightarrow\int u\MA(v),&
\int b_j\MA(b_j)&\longrightarrow\int v\MA(v).
\end{align*}
Consequently equal actual measures $\MA(u)=\MA(v)$ imply
\begin{equation}\label{403Ilimit}
I_j=\int(a_j-b_j)[\MA(b_j)-\MA(a_j)]\longrightarrow0.
\end{equation}
All these integrals are finite by the proved full-dual bound.

\section{Every bounded test gradient energy tends to zero}
Subtract fixed upper constants from $a_j,b_j$ and any fixed
$w\in\C$, and choose $C\ge0$ so their energies are at least
$-C$. The normalized functions are nonpositive. Fixed shifts
do not change curvatures, gradient pairings or $I_j$, because
all measures have mass $M$.
For any such $p,q$, equation~\eqref{403cocycle} at zero gives
$-\int q\MA(q)\le3C$, whence $E^*(\MA(q))\le2C$ and
\[
\int(-p)\MA(q)\le3C.
\]
Also $B(q,r)\le2\MA((q+r)/2)$, and concavity bounds the
midpoint energy below by $-C$. Hence
$\int(-p)B(q,r)\le6C$.

Write $f=a-b$, $g=(a+b)/2$, $h=w-g$. Green and symmetry give
\begin{align*}
I(a,b)&=Q_{\theta_a}(f,f)+Q_{\theta_b}(f,f),\\
Q_{\theta_w}(f,f)-I(a,b)/2
&=Q_{\theta_a}(f,h)-Q_{\theta_b}(f,h).
\end{align*}
The second identity has both sides equal to
$-\int f\DD f\wedge\DD h$. Since
$|h|\le-w+(-a-b)/2$, the preceding moment bounds give
$\int|h|\MA(q)\le6C$ and $\int|h|B(q,r)\le12C$.
Expanding
\[
Q_{\theta_a}(h,h)=
-\int h[B(a,w)-(\MA(a)+B(a,b))/2]
\]
bounds it by $12C+3C+6C=21C$, and the same holds for the
$b$ slot. Cauchy--Schwarz in the direct sum of the two positive
Hilbert spaces yields
\begin{equation}\label{403alltest}
0\le Q_{\theta_w}(a-b,a-b)
\le I(a,b)/2+\sqrt{42C I(a,b)}.
\end{equation}
From~\eqref{403Ilimit}, for $d_j=a_j-b_j$,
\begin{equation}\label{403Qlimit}
Q_{\theta_w}(d_j,d_j)\longrightarrow0
\quad\text{for every fixed }w\in\C.
\end{equation}
There is no uniform sup-norm assumption on the presentations.

\section{Finite Gauss-disc lifts and a curve endpoint inequality}
For a complete nontrivially valued field $F$ containing $K$,
write $d_F(q)=\operatorname{trdeg}_{\widetilde F}
\widetilde{\HH(q)}$ in the graded sense.

\begin{lemma}\label{403disc}
If $q\in(\mathbb G_m)^{\rm an}_F$ has $d_F(q)=1$ and
$|z(q)|=r$, some finite extension $E/F$ has a lift
$q_E=\eta_{a,R}$ with $a\in E$, $0<R\le r$.
The path $\eta_{a,t}$, $R\le t\le r$, ends at
$\eta_{0,r}$, has finite log-length $\ell=\log(r/R)$,
and has $|z|=r$ throughout.
\end{lemma}
\begin{proof}
Let $C_F$ be the completion of an algebraic closure of $F$.
Choose a lift of $q$ to this scalar extension. Existence is the
standard surjectivity of analytic scalar extension. The full
graded-dimensional line point is an Abhyankar presentation, so
CLD 13.8.8(d) says its lifted graded residue field is finite over
the compositum of the old one and $\widetilde{C_F}$.
Completion adds no initial forms, and $\widetilde{C_F}$ is
algebraic over $\widetilde F$. Therefore the lift still has
graded dimension one.

Here is the classification needed over the algebraically closed
field $C_F$, proved from the norm. For its lifted seminorm let
$R=\inf_a|z-a|(q)$, and choose $a_n$ with
$R_n=|z-a_n|(q)\downarrow R$. The closed discs
$D(a_n,R_n)$ are nested, since
$|a_n-a_m|\le\max(R_n,R_m)$. If the infimum is attained,
use that center directly. Otherwise this gives a strictly
decreasing cofinal sequence of radii.
If the discs have an intersection point $a$, then
$|z-a|(q)=R$. For every $b$, the ultrametric inequality and
the definition of $R$ give
$|z-b|(q)=\max(R,|a-b|)$. Factoring every polynomial into
linear factors shows the seminorm is exactly $\eta_{a,R}$.
If $R=0$ it is rigid and has graded dimension zero.
If the intersection is empty, every $b$ eventually lies
outside a disc, so $|a_n-b|>R_n$ and the initial form of
$z-b$ is that of the base constant $a_n-b$.
Every nonzero polynomial and rational function then has initial
form in the base graded field. Completion again adds no initial
forms. The resulting graded dimension is zero, a contradiction.
Thus $R>0$ and the lift is a genuine positive-radius Gauss disc.

Approximate $a$ by an algebraic $a_0$ with $|a-a_0|<R$.
The Gauss discs with centers $a,a_0$ coincide. Set
$E=F(a_0)$; the point $\eta_{a_0,R}$ over $E$ projects to $q$.
Its norm of $z$ is $r=\max(|a_0|,R)$. Thus $R\le r$;
if $R<r$, $|a_0|=r$, and if $R=r$, $|a_0|\le r$.
The stated path stays in $\mathbb G_m$, has norm of $z$
equal to $r$, and at its last disc equals $\eta_{0,r}$.
\end{proof}

We spell out the smooth curve energy normalization. On a Gauss
interval $\eta_{a,t}$ in a smooth affine line, use $z-a$ as
a moment and set $x=\log t$. Refine any finite set of smooth
moment charts by this coordinate. On a compact interval their
analytic-unit log norms are piecewise affine in $x$, with
integer slopes and finitely many breaks: the Gauss norm is the
maximum of Laurent monomial norms, and a compact interval
inside the analytic domain meets finitely many slopes. The
projection to $z-a$ is an analytic open embedding near each
Gauss point. Its local finite flat degree is one. CLD 14.5.2--3
therefore gives canonical calibration one on this interval;
the moment vector's first coordinate has slope one and is
primitive. For a smooth function $h$, its restriction $H(x)$
is continuous and piecewise smooth. The positive curve energy
on the interval is exactly $|H'(x)|^2\,dx$. Other branches
have nonnegative energy. Break points contribute no negative
terms to an integral of smooth $(1,1)$ forms. Consequently,
if $\chi\ge0$ is a smooth cutoff equal to one on the path,
\begin{equation}\label{403trace}
|h(\eta_{a,R})-h(\eta_{0,r})|^2
\le\ell\int\chi\,d'h\wedge d''h.
\end{equation}
This is ordinary Cauchy--Schwarz on the finite real interval.
The case $\ell=0$ has identical endpoints.

For a finite scalar extension $E/F$, the normalized
varying-ground-field identity is
\begin{equation}\label{403base}
\int_{C_E}p^*(\chi\,d'h\wedge d''h)
=\int_C\chi\,d'h\wedge d''h.
\end{equation}
It is precisely CLD 15.3.1, equivalently the fundamental
current identity 16.1.1. The support is compact because the
finite scalar-extension map is proper. There is no factor
$[E:F]$ in this operation. The extension to uniform/strong
gradient limits follows by applying it to smooth differences
and taking their Hilbert norms.

In applications $h$ is an ambient continuous dpsh function.
Its local uniform smooth approximants with strong gradients
are supplied by the appendix, for the chosen positive test
current. The next section bounds the cutoff curve norms of
their differences by those ambient norms. Equation~\eqref{403base}
gives the same control after the finite extension in
Lemma~\ref{403disc}. The restrictions converge uniformly,
including at both endpoints, and strongly in cutoff curve
energy. Applying~\eqref{403trace} to the smooth restrictions
and passing both limits proves the same endpoint inequality
for this continuous restriction. No unbounded limit function
is pulled back before this step.

\section{Positive actual Gauss fibres and their global bounded caps}
Let $P(z_2)$ be a nonconstant one-variable Laurent polynomial,
let $r_P>0$, and put
\[
\psi=\max(0,\log(|P(z_2)|/r_P)),\qquad T=\DD\psi.
\]
CLD 18.6.2 with $d=1$ says exactly
\begin{equation}\label{403fibre}
T=\delta_{P^{-1}(\eta_{r_P})}
\end{equation}
on the boundaryless torus surface. The fibre is viewed over
$F_0=\HH(\eta_{r_P})$. Its equation is $P(z_2)=T_0$,
where $T_0$ is the transcendental coordinate at the Gauss
point, and $z_1$ is free. The finite Laurent equation has
no zero root and is squarefree in characteristic zero.
Indeed a repeated root would make the transcendental $T_0$
a critical value algebraic over the ground field. Thus the
fibre is a finite disjoint union of smooth curves
$C_\nu=(\mathbb G_m)^{\rm an}_{L_\nu}$ with
$[L_\nu:F_0]<\infty$. When these curves are viewed over
$F_0$, their integration calibration is
$\kappa_\nu=[L_\nu:F_0]>0$ times the normalization over
$L_\nu$, by CLD 14.4.3. This is distinct from
equation~\eqref{403base}; there is no assertion that
$L_\nu/K$ is finite.

Fix a small injective tropical box around the entire path
in question, and use the centered reference
$\rho_a=|\tau-a|^2/2$. Choose $\delta>0$ small relative to
the box width and $\epsilon>0$ small enough that
\[
w=\max(\rho_a,\delta+\epsilon\psi)-\rho_a
\]
is zero on a fixed tropical boundary collar. The function
$\psi$ is nonnegative and bounded above on the compact
polyannular box; the Laurent polynomial bound gives the
required choice of $\epsilon$. Extend $w$ by zero to $X$.
It belongs to $\C$. To check the zero locus of $P$ and all
continuous stages explicitly, take a smooth convex $\chi_0$
equal to zero below zero and to $x-c$ above one. For $N>3$,
\[
p_N=\chi_0(\log(|P|/r_P)+N)-N+c
\]
is finite smooth psh, including near zeros, and
$\psi=\max(0,p_N)$ exactly. Smooth convex regularizations
of this inner maximum and of the outer maximum converge
uniformly and preserve a fixed boundary gap. Their zero
extensions are global smooth $\theta$-psh functions.
This proves the asserted membership without a first-stage
identification of the full sheaf.

On the fixed inner region $\rho_a<\delta/2$, the outer
maximum chooses its second slot, so
$\theta_w=\epsilon T$. Let a nonnegative smooth cutoff
$\chi$ be supported there and equal to one on the whole
path. For smooth ambient $h$, positivity and
equation~\eqref{403fibre} give
\begin{equation}\label{403currentbound}
Q_{\theta_w}(h,h)
\ge\epsilon\kappa_\nu
\int_{C_\nu}\chi\,d'(h|_{C_\nu})\wedge d''(h|_{C_\nu}).
\end{equation}
Smooth restriction here is ordinary analytic pullback of
moments, with fibre base constants constant on the curve.
The integration current is precisely the source theorem's
fibre current; no pullback of an arbitrary positive current
is postulated. Uniform/strong ambient approximation extends
\eqref{403currentbound} to continuous dpsh $h$, as described
after~\eqref{403base}. In detail, the right hand norm of
every difference of approximants is bounded by the left
hand norm on the fixed inner cutoff. Hence the restrictions
are strongly Cauchy; uniform convergence identifies their
function limits and gradient uniqueness identifies their
curve gradients. The positive bound survives the limit.

Apply a finite disc lift to a full-dimensional point $q$
of one fibre curve, and let $q_0$ be its radius-$r$ Gauss
endpoint. If the projected whole path lies in the inner
region, equations~\eqref{403trace}, \eqref{403base} and
\eqref{403currentbound} give, for every ambient continuous
dpsh $h$,
\begin{equation}\label{403endpoint}
|h(q)-h(q_0)|^2
\le\frac{\ell}{\epsilon\kappa_\nu}
Q_{\theta_w}(h,h).
\end{equation}
Every coefficient is fixed and finite. The lift's scalar
degree is not inserted in $\kappa_\nu$.

\section{Two actual coordinate fibres connect every full point}
Let $x\in Y$ have $d_K(x)=2$ and write
$r_i=|z_i(x)|$. The projection $q_2$ to the second
coordinate has graded dimension one. In fact its graded
dimension is at most one, the relative fibre is a line
over $\HH(q_2)$ with relative graded dimension at most one,
and graded transcendence degrees add in the residue-field
tower. Total dimension two forces both dimensions to be
one.

The support of $q_2$ on $K[z_2^{\pm1}]$ is zero.
Completion of its fraction field adds no initial forms,
so the homogeneous initials of all Laurent polynomials
generate its graded residue field by fractions. Choose a
Laurent polynomial $P(z_2)$ with transcendental initial
form. No highest-norm terms in a polynomial in $P$ can
cancel. Therefore $P(q_2)=\eta_{r_P}$ with
$r_P=|P(q_2)|>0$. The point $q_2$ is one branch of the
finite fibre of $P$ at this Gauss point, with
$L=\HH(q_2)$ finite over $\HH(\eta_{r_P})$.

The corresponding actual surface fibre component is
$C=(\mathbb G_m)^{\rm an}_L$ with variable $z_1$.
The point $x$ is full graded-dimensional on this curve
by the same tower calculation. Lemma~\ref{403disc} gives
a finite lifted path from $x$ to
\[
\xi_1=\eta_{L,r_1}\quad\text{over }q_2.
\]
On the whole path both original coordinate norms are
fixed at $(r_1,r_2)$. It therefore lies in the inner
tropical region of one bounded cap of $P(z_2)$.
Equation~\eqref{403endpoint} and
\eqref{403Qlimit}, applied to $h=d_j$, prove
\begin{equation}\label{403first}
d_j(x)-d_j(\xi_1)\longrightarrow0.
\end{equation}
The point $\xi_1$ has total graded dimension two: its
new independent Gauss $z_1$ lies over the graded-dimension
one field $L/K$. Its projection in the first coordinate
is $\eta_{K,r_1}$.

Now regard $\xi_1$ as a point of the second actual fibre
\[
C'=z_1^{-1}(\eta_{K,r_1})
=(\mathbb G_m)^{\rm an}_{\HH(\eta_{K,r_1})}
\]
with variable $z_2$. Its relative graded dimension is
one, again by additivity. Lemma~\ref{403disc} connects
it, after a finite extension of this curve's own field,
to the joint Gauss point
$\xi=\eta_{K,(r_1,r_2)}$. The two original norms are
again constant on the path. Use the bounded cap of
$\max(0,\log(|z_1|/r_1))$, whose current is the
integration current of $C'$, to obtain
\begin{equation}\label{403second}
d_j(\xi_1)-d_j(\xi)\longrightarrow0.
\end{equation}
These are two genuine analytic fibres. There is no
claim that the Shilov section over a varying base point
is analytic; that incorrect shortcut would contradict
the explicit warning in Gubler--Rabinoff 5.7.

The appendix supplies, at every total graded-dimension
two point $y$, a bounded $w_y\in\C$ with
$\MA(w_y)\ge\alpha_y\delta_y$, $\alpha_y>0$.
The finite full dual then shows $u(y)$ and $v(y)$ are
finite. Consequently the decreasing presentations have
finite endpoint limits at $x,\xi_1,\xi$. Passing those
limits in~\eqref{403first}--\eqref{403second} gives
\begin{equation}\label{403fulltoskeleton}
(u-v)(x)=(u-v)(\eta_{K,(r_1,r_2)}).
\end{equation}

\section{One skeleton constant and pointwise Shilov reconstruction}
Fix $r_2$ and vary $\log r_1$ along a finite interval.
The resulting joint Gauss points lie on the actual
curve $z_2^{-1}(\eta_{K,r_2})$ over
$\HH(\eta_{K,r_2})$. Partition the interval into
finitely many pieces small enough to fit in inner
injective tropical boxes. A cap of $z_2$ provides
the positive curve current on each piece.
For a smooth function the direct Gauss-radius energy
calculation gives the endpoint inequality with length
the piece's log-radius length; the same strong-limit
argument extends it to $d_j$. Equation~\eqref{403Qlimit}
therefore makes $u-v$ equal at the two endpoints of
every piece. The endpoint values are finite by the
point atoms. Varying $r_2$ at fixed $r_1$ uses the
other coordinate curve in exactly this way.
Horizontal and vertical finite log-radius paths
connect any two points of $\RR^2$ before passage to
the period quotient. Thus $u-v$ is one finite constant
$c$ on the entire reference Gauss section. Equation
\eqref{403fulltoskeleton} proves the same equality at
every point of total graded dimension two on $X$.

For every sufficiently small affinoid neighbourhood
$A$ in a local weight chart, all its Shilov boundary
points have graded dimension two. This is the precise
application of Gubler--Rabinoff Lemma 6.3(3)--(4) to
our smooth pure two-dimensional space, followed by
the Abhyankar upper bound. All-stage continuous CLD
recognition and decreasing closure make $\rho+u$
and $\rho+v$ classically psh, including at poles.
The appendix proves this recognition by the
uniform-to-decreasing construction and states the
source's complete-field and curve hypotheses.
The USC Shilov maximum principle therefore gives
\[
\max_A(\rho+u)=\max_{\Gamma_A}(\rho+u)
=\max_{\Gamma_A}(\rho+v)+c=\max_A(\rho+v)+c.
\]
For any point $p$, take the infimum over arbitrarily
small affinoid neighbourhoods. USC and the maximum
principle give exactly
$q(p)=\inf_{A\ni p}\max_{\Gamma_A}q$ for each local
weight $q$, including $q(p)=-\infty$.
It follows that $\rho+u=\rho+v+c$ pointwise,
and hence $u=v+c$ on all of $X$.

\begin{theorem}\label{403uniqueness}
For actual $u,v\in\G$, equality of the actual
Monge--Ampere measures implies $u=v+c$ everywhere
for one finite real constant $c$. If both are
normalized by $\sup_Xu=\sup_Xv=0$, they coincide.
\end{theorem}
\begin{proof}
The all-four-moment limit proves~\eqref{403Ilimit};
the all-test estimate proves~\eqref{403Qlimit};
the finite curve inequalities, point atoms, two
actual fibres and the preceding Shilov argument
prove the claimed pointwise conclusion. Generated
weights are bounded above and not identically
$-\infty$, so their suprema are finite. Equal zero
suprema force $c=0$.
\end{proof}

\section{Scope and residual original requirements}
This theorem concerns the actual decreasing-generated
class only. It does not identify that class with a
maximal intrinsic $\mathcal E^1$, construct an actual
generated limit for every normalized energy-bounded
family, prove an envelope theorem, or solve the
Monge--Ampere equation for every prescribed finite-dual
measure. In particular product compactness of the
Shilov values alone does not prove representability.
Countable energy-bounded decreasing closure from
paper 400 remains a separate proved working claim.
No original problem count or acceptance count changes.

All mathematical assertions new in this paper are
UNREVIEWED. The point-atom appendix is a byte-identical
mathematical-body reproduction of paper 402, not a
second independent source audit. The original source
bodies, literal units, pixels and book contexts remain
local only. Source reading verifies the stated
application boundary, not correctness by hash. Native
and supplied programs are data and were not executed
or imported. No active native campaign has received
this paper as hot feedback.

\clearpage
\appendix
\section*{Whole mathematical body of point-atom working paper 402}
The following is the exact mathematical body of paper 402,
including its historical scope and residual statements. Those
statements describe what paper 402 alone establishes; the new
generated uniqueness argument is the main part of this paper.
Its inclusion supplies the full bounded Green, local atomic,
all-stage recognition and pointwise reconstruction proofs rather
than an external project-file dependency. Exact identity is not
an additional independent primary-source reading.
\section{The actual analytic object and proof foundations}
Put $K=\mathbb C((s))$, $|s|=e^{-1/8}$, $t=s^8$,
$Y=(\mathbb G_m^2)^{\mathrm{an}}_K$, and
\[
 X=Y/\langle(t,2),(3,t)\rangle,\qquad
 \tau=(-\log|z_1|,-\log|z_2|).
\]
The translations on $\tau$ form $\mathbb Z^2$. Ordinary open tropical
boxes of width less than one give injective quotient charts. We use
the previously reconstructed analytic quotient: compact, separated,
good, boundaryless, connected, smooth, and pure of dimension two.
The quotient construction and its nonalgebraization proof are not
reproved here and are not consequences of the present point test.
On a chart let $\rho=|\tau|^2/2$ and $\theta=\DD\rho$.
Changes of reference are affine unit logarithms. With the identity
torus's saturated lattice and degree one,
\[
 M=\int_X\theta^2=2.
\]
This is the raw CLD convention; no probability normalization is made.

Let $\mathscr P$ be the minimal continuous CLD psh sheaf containing
the ordinary smooth psh functions and closed under local uniform
limits. All stages of this sheaf are retained. Set
\[
 \C=\{a\in C(X,\RR):\rho+a\in\mathscr P\text{ locally}\},\quad
 \theta_a=\theta+\DD a,\quad B(a,b)=\theta_a\wedge\theta_b.
\]
We use the actual recursive continuous products. Their positivity,
symmetry, closedness, locality, local uniform continuity, and
invariance under smooth $\DD$-null reference changes are precisely
CLDv2 8.5.1, 8.5.3, 8.5.5--6, 8.7.1--4 and 8.8.4--5, at the stated
continuous psh/dpsh and positive closed-current hypotheses. Positive
top currents are Radon measures by 7.2.8 on these locally
differentially separated charts. Smooth cutoffs, partitions, scalar
approximation and boundaryless Stokes are the conventional analytic
foundations. Arbitrary continuous-multiplier closedness in 8.6.7 is
not a premise. No first-stage identity for $\mathscr P$ is assumed.

Constants, convex combinations and finite maxima preserve $\C$.
For a smooth slot, convex regularized maxima are smooth psh.
Uniform continuity and sheaf minimality extend the operation one
slot at a time to all stages. This also proves ordinary convex
composition with finite continuous psh slots by the same argument.

\section{Bounded Green calculus and the full bounded dual}
Here we give the extra energy argument rather than infer it from
a positivity label. For a positive closed $(1,1)$ current $T$,
complete compact smooth $(1,0)$ forms modulo the null space of
\[
 (a,b)\longmapsto\int a\wedge Jb\wedge T.
\]
Use a nonnegative compact smooth cutoff $\chi$ for the corresponding
local seminorm $Q_{T,\chi}$. Let $\HH_T$ consist of continuous
functions locally uniformly approximable by smooth functions whose
gradients converge strongly in these seminorms. The gradient is
unique: for a uniformly zero sequence $f_j$ with a strong gradient
limit and a compact smooth form $b$, closedness gives
\[
 \int d'f_j\wedge Jb\wedge T=-\int f_jd'(Jb)\wedge T\longrightarrow0.
\]
The right-hand coefficient measure is finite on the fixed compact
support; smooth Hilbert tests are dense. The limit gradient is zero.
Partitions patch local approximants by the Leibniz rule.
Bounded continuous coefficients act by scalar approximation and the
bound $\|fa\|\le\|f\|_\infty\|a\|$. Products and cutoffs have their
usual gradient rules.

If $f\in\HH_T$ and $L_Tf=\DD(fT)$ is signed Radon, smooth
approximation first in $f$, then uniformly and strongly in a
compactly supported $g$, yields
\[
 \int gL_Tf=-Q_T(f,g),\qquad
 Q_{T,\chi}(w,w)=\frac12\int w^2\DD\chi\wedge T-\int\chi wL_Tw.
\]
The sheaf of $f\in\HH_T$ with $L_Tf\ge0$ contains smooth psh
functions and is uniformly closed. Indeed take nested cutoffs
$\chi,\eta$, $\eta=1$ near $\operatorname{supp}\chi$,
a common bound $|f_j|\le B$ on the whole support of $\eta$, and
$\epsilon=\|f_j-f_k\|$ there. Positivity and the preceding identity give
\[
 0\le Q_{T,\chi}(f_j-f_k,f_j-f_k)
 \le\frac{\epsilon^2}{2}\|\DD\chi\wedge T\|_{\mathrm{TV}}
 +2B\epsilon\|\chi\|_\infty\|\DD\eta\wedge T\|_{\mathrm{TV}}.
\]
Choose smooth approximants to $f_j$ with uniform and smaller-patch
gradient errors below $1/j$. The estimate proves strong convergence;
distributional positivity survives. Minimality places every stage
of $\mathscr P$ in this sheaf. Its vector span includes continuous
dpsh functions. Finite patching on compact $X$ proves
\[
 \int f\DD(gT)=\int g\DD(fT)=-Q_T(f,g),\qquad Q_T(f,f)\ge0.
\]
In particular all the following computations hold at every
continuous stage. The exact terms in the bounded product expansion
have zero total integral, so every $B(a,b)$ has mass $M=2$.

Define
\[
 E(a)=\frac13\int_X a[\theta^2+\theta_a\wedge\theta+\theta_a^2].
\]
For $a_t=(1-t)a+tb$, the cubic expansion and Green give
\[
 \frac{d}{dt}E(a_t)=\int(b-a)\MA(a_t),\qquad
 \frac{d^2}{dt^2}E(a_t)=-2Q_{\theta_{a_t}}(b-a,b-a)\le0.
\]
Thus $E$ is concave and monotone, $E(0)=0$, and
$E(a+c)=E(a)+Mc$. Its tangent inequality is
\[
 E(a)\le E(w)+\int(a-w)\MA(w)\qquad(a,w\in\C).
\]
For a positive Radon measure $\mu$ define the \emph{full} bounded dual
\[
 E^*(\mu)=\sup_{a\in\C}\left(E(a)-\int a\,d\mu\right).
\]
Constants force $\mu(X)=M$ whenever this quantity is finite.
The tangent inequality, with equality at $a=w$, proves the exact
formula
\begin{equation}\label{dual}
 E^*(\MA(w))=E(w)-\int w\MA(w)<\infty
 \quad\text{for every }w\in\C.
\end{equation}
It is nonnegative, since $a=0$ is one of the competitors.

\section{The actual generated class and a moment bound}
Let $\G$ consist of actual pointwise limits
$a_j\downarrow u$ on $X$, where $a_j\in\C$ and
$\inf_jE(a_j)>-\infty$. These limits are USC.
They are not abstract functionals or equivalence classes modulo
null sets. Define $E(u)=\lim_jE(a_j)$.
This is independent of every continuous decreasing presentation.
To see this, if $b_k\downarrow u$ and $a$ is continuous with $u<a$
everywhere, the increasing ordinary open sets $\{b_k<a\}$ cover
compact $X$. One index therefore works everywhere. Apply this to
$a=a_j+\epsilon$, use monotonicity and the constant rule, and
then reverse the roles of the presentations. The limits agree.
In particular $E(u)>-\infty$ for every presentation of a member
of $\G$. The identically $-\infty$ limit is excluded: compact
strict domination would eventually put a presentation below every
negative constant, contradicting its energy lower bound.

\begin{lemma}\label{moment}
If $u\in\G$, $u\le0$, and $E^*(\mu)<\infty$ for a positive
Radon measure $\mu$, then
\[
 \int(-u)\,d\mu\le E^*(\mu)-E(u)<\infty.
\]
\end{lemma}
\begin{proof}
Every term of a presentation satisfies
$\int(-a_j)\,d\mu\le E^*(\mu)-E(a_j)$.
The continuous first term bounds all terms above. Apply monotone
convergence to that common bound minus $a_j$ and pass to the limit.
This proves the claim even at poles and does not presume
integrability before the bound. The same argument proves
$L^1(\mu)$ convergence of the whole presentation.
\end{proof}

\section{Laurent Abhyankar moments without changing the field}
Write $d_K(x)=\operatorname{trdeg}_{\widetilde K}
\widetilde{\mathscr H(x)}$, with graded residue fields.
Let $\mathscr A$ be the set of points of $X$ with $d_K(x)=2$.
A lift of such a point in $Y$ will also be denoted $x$.

\begin{lemma}\label{moments}
There are two nonzero Laurent polynomials $P_1,P_2$ over $K$
whose initial forms at $x$ are algebraically independent over
$\widetilde K$. Put $r_i=|P_i(x)|>0$. Then
\[
 (P_1,P_2)(x)=\eta_{(r_1,r_2)}
\]
in $(\mathbb A_K^2)^{\mathrm{an}}$. The map is finite and flat at
$x$, and admits a local finite flat representative with a positive
constant degree $D$ and with $x$ the only point above this Gauss point.
\end{lemma}
\begin{proof}
The support of $x$ on the Laurent polynomial ring is zero.
Otherwise it lies on a proper algebraic subvariety of dimension
at most one, contradicting $d_K(x)=2$ and the Abhyankar inequality.
Thus $\mathscr H(x)$ is the completion of the valued field
$K(z_1,z_2)$. Completion does not add graded residue elements:
approximate a nonzero element with error strictly less than its
norm. Its initial form agrees with that of the approximant.
Consequently the initial forms of all nonzero Laurent polynomials
generate the graded residue \emph{field}, allowing ratios.
A maximal algebraically independent subfamily of these homogeneous
generators has cardinal two. Choosing its members gives $P_1,P_2$.

For every polynomial $\sum_I c_IP^I$, the initial terms of maximal
norm cannot cancel by graded algebraic independence. Its norm is
$\max_I|c_I|r^I$. This is exactly the Gauss seminorm, for arbitrary
real $r_i>0$; no scaling of $r_i$ into $|K^\times|$ is required.
This is also the Abhyankar-basis criterion of CLD 13.8.5.

The torus is boundaryless, hence $x$ is in the relative interior
of the map to $\mathbb A^2$. The target is reduced and the Gauss
point has graded dimension two. CLD 11.6.5(c) makes the map finite
and flat at $x$. Its 11.6.3 local representative can be chosen
with a single point of that fibre. The rank is a positive integer
$D$. After shrinking the target inside its ordinary interior,
the source is an ordinary open neighborhood of $x$, and the finite
flat degree remains $D$ there.
\end{proof}

\section{The actual atomic mass and its global test}
We spell out the local calibrated number.
On a real two-plane set $g(t)=\max(0,t_1,t_2)$ and
$g_\epsilon=\epsilon\log(1+e^{t_1/\epsilon}+e^{t_2/\epsilon})$.
For $g_1$, write $p_i=e^{t_i}/(1+e^{t_1}+e^{t_2})$ and
$p_0=1-p_1-p_2$. The Hessian determinant is $p_0p_1p_2$.
The gradient maps the plane diffeomorphically to
$\{p_1,p_2>0:p_1+p_2<1\}$ with that Jacobian.
The raw superform rule is $2!\det(\operatorname{Hess}g_1)\,dt_1dt_2$.
Its integral is $2!\,(1/2)=1$.
Dilation and dominated convergence concentrate this finite measure
at the origin as $\epsilon\downarrow0$, so
\[
 (\DD g)^2=\delta_0.
\]
This is the explicit case of CLD 1.7.1 and 1.7.6--8, including
their factor $2!$.

Near $x$, both $P_i$ are invertible. Apply the pleasant moment
$(P_1,P_2)$ there. Finite flatness and the single Gauss fibre imply
full tropical dimension and that the calibrated tropical cone at
$x$ is the whole real plane with multiplicity $D$.
Indeed the map is open, so the Gauss section of a neighborhood of
the target gives a neighborhood in the tropical plane. The explicit
calibration formula CLD 14.5.2--3 makes its image calibration
$D$ times the standard one. Locality removes other fibre points,
since there are none in this representative. More explicitly,
smooth log-sum-exp approximants have top measures strongly
supported on the inverse image of the target Gauss section.
The calibrated integration formula pushes them to $D$ times the
corresponding real-plane measures near the target tropical value.
Those measures concentrate there with mass one. Finite properness
and the single Gauss fibre imply that every inverse Gauss point
over a sufficiently small tropical neighborhood lies in any
prescribed ordinary neighborhood of $x$: otherwise compactness
would give another point of that fibre. Continuous tests therefore
converge to $D$ times evaluation at $x$. Local uniform continuity
of the actual product identifies the limit with the log-maximum
product. This reconstructs the atomic calculation 8.9.2--4 and yields
\begin{equation}\label{localatom}
 \left(\DD\max(0,\log(|P_1|/r_1),\log(|P_2|/r_2))\right)^2
 =D\delta_x
 \quad\text{on some ordinary neighborhood of }x.
\end{equation}
Formula \eqref{localatom} uses the normalized
canonical calibration over the unchanged field $K$ and the actual
finite flat degree. No arbitrary-current pullback is invoked.
CLD 18.6.2 gives the equivalent Gauss-fibre description, but its
relative-intersection proof is not needed for this argument.

\begin{theorem}\label{atom}
For every $x\in\mathscr A$ there is $w_x\in\C$ and
$\alpha_x>0$ such that the actual positive Radon measure
$\MA(w_x)$ satisfies
\[
 \MA(w_x)(\{x\})=\alpha_x.
\]
One can take $\alpha_x=\varepsilon^2D$ in the construction below.
\end{theorem}
\begin{proof}
Choose an injective lift chart around the tropical class
$a=\tau(x)$ and a small Euclidean tropical ball of radius $R$ there.
Its inverse image has compact closure inside a polyannulus.
Set $\rho_a=|\tau-a|^2/2$; its difference from $\rho$ is
smooth affine and $\DD$-null.
The function
\[
 \psi=\max(0,\log(|P_1|/r_1),\log(|P_2|/r_2))
\]
is finite, continuous, nonnegative and psh on this whole chart,
including zeros of the polynomials. Here is a direct check at all
continuous stages. Choose a smooth convex $\chi$ equal to zero for
$v\le0$, with $0\le\chi'\le1$, and equal to $v-c$ for $v\ge1$,
where $0<c<1$. For $N>3$ the finite smooth psh function
\[
 p_{i,N}=\chi(\log(|P_i|/r_i)+N)-N+c
\]
is locally constant at zeros. Where it can be positive it equals
$\log(|P_i|/r_i)$, so $\psi=\max(0,p_{1,N},p_{2,N})$ exactly.
Smooth convex regularized maxima place it in $\mathscr P$.

Take $0<\delta<R^2/4$ and choose $\varepsilon>0$ so small that
$\delta+\varepsilon\sup\psi<R^2/2$ on the compact closure.
Then
\[
 w_x=\max(\rho_a,\delta+\varepsilon\psi)-\rho_a
\]
vanishes on a tropical boundary collar. Extend it by zero to $X$.
It belongs to $\C$: regularize the inner finite maximum, and then
use a smooth convex regularized outer maximum which agrees exactly
with the larger slot when the gap exceeds its regularization
width. The fixed collar gap ensures that each sufficiently small
regularization extends by zero smoothly and $\theta$-psh.
Their uniform limit is the displayed function.

On the fixed inner neighborhood $\rho_a<\delta/2$ we have
$\theta_{w_x}=\varepsilon\DD\psi$. Shrink further to the
neighborhood in \eqref{localatom}. Locality of the actual recursive
product gives
$\MA(w_x)=\varepsilon^2D\delta_x$ there.
Outside it the global product remains positive and has raw mass
$M=2$. Thus its atom at $x$ is precisely the positive number
$\alpha_x=\varepsilon^2D$.
\end{proof}

\section{Finiteness, sublevel bounds, and the precise compactness limit}
\begin{corollary}\label{finite}
For every $u\in\G$ and every $x\in\mathscr A$, $u(x)$ is finite.
More precisely, if $\sup_Xu=0$ and $E(u)\ge-C$, then
\[
 -\frac{E^*(\MA(w_x))+C}{\alpha_x}\le u(x)\le0 .
\]
\end{corollary}
\begin{proof}
Use \eqref{dual} and Lemma \ref{moment} for the \emph{whole}
positive measure $\mu_x=\MA(w_x)$.
Since $-u\ge0$ and $\mu_x$ contains the positive atom
$\alpha_x\delta_x$,
\[
 \alpha_x(-u(x))\le\int(-u)\,d\mu_x
 \le E^*(\mu_x)-E(u)\le E^*(\mu_x)+C.
\]
This excludes a pole at $x$ and proves the bound.
An arbitrary $u$ can be shifted by its finite supremum.
A member of $\G$ has a finite upper bound and is not identically
$-\infty$, so that shift is legitimate.
\end{proof}

Every Shilov point of a pure-dimensional two-dimensional affinoid
domain in a local chart has $d_K=2$. At the exact general-radius
hypotheses, GRv2 Lemma6.3(3)--(4) expresses these Shilov points as
the strata with two graded algebraically independent generators;
the Abhyankar inequality gives the reverse bound.
The entire lemma and its immediate proof were read, including its
Temkin graded-reduction and Ducros fibre-dimension foundations.
These deeper theorems are named published foundations, not freshly
recursively reproved here. Hence Corollary \ref{finite} applies to
every such Shilov point, not only to the reference torus skeleton.
On a fixed finite Shilov boundary the lower bound can be chosen
uniformly by taking the maximum of the finitely many displayed
constants.

For clarity, actual local weights are classically psh at all
stages. The published smooth CLD-to-classical bridge is the
unnumbered statement and argument in GRv2 Section1.2.
Its Proposition4.2 gives sheafiness and decreasing-net
closure. A uniform limit of finite continuous classical psh
functions is classical psh: select errors at most $2^{-j}$ and add
$3\cdot2^{-j}$, obtaining a decreasing sequence with the same limit.
Sheaf minimality then recognizes all of $\mathscr P$, and a
continuous decreasing presentation recognizes $\rho+u$ at poles.
GRv2 Theorem6.4 therefore gives the affinoid Shilov maximum
principle for these actual USC weights.
For all arbitrarily small affinoid neighborhoods $A$ in one chart,
\[
 (\rho+u)(y)=\inf_{A\ni y\ \mathrm{neighborhood}}
                \max_{\xi\in\Gamma_A}(\rho+u)(\xi).
\]
The maximum principle gives one inequality. USC gives the other
by every finite upper bound at a pole and by
$(\rho+u)(y)+\eta$ at a finite value. Thus values on the union
$\mathscr S$ of \emph{all} these Shilov boundaries determine the
actual weight, including all poles.

Fix a normalized energy sublevel
$\G_C=\{u\in\G:\sup u=0,\ E(u)\ge-C\}$.
Its evaluation map on $\mathscr S$ is injective and takes values in
\[
 \prod_{\xi\in\mathscr S}[-L_\xi(C),0],\qquad
 L_\xi(C)=\frac{E^*(\MA(w_\xi))+C}{\alpha_\xi}.
\]
Tychonoff compactness implies that every net of such evaluations
has a coordinatewise convergent subnet with finite coordinates.
This is compactness in the product space of boundary values.
It does \emph{not} show that the limit coordinates are represented
by an actual member of $\G$, that the normalized image is closed,
or that energy is semicontinuous there. Those are still mathematical
gates for actual energy-sublevel compactness.
Likewise this point bound alone does not prove uniqueness from
equality of generated MA measures.

\begin{remark}
We have not claimed that $\delta_x$ itself has finite full bounded
dual merely because it is part of $\MA(w_x)$.
The whole test measure has the supporting-energy formula
\eqref{dual}; its positivity is exactly what gives the atomic
point inequality. The distinction retains the raw mass-two dual
and avoids an invalid componentwise assertion.
\end{remark}

\section{Source scope and original residual}
The exact CLDv2 full source is bound locally by SHA256
\texttt{64846cb00926267c6238d1afac61b5eaed0a0e3a6349725738f9b1bf163de085}.
Applied immediate readings include 1.7.1 and1.7.6--8
(LF1773--1830 and1931--2135), the complete local finite-flat block
11.6.1--5 (LF15846--15912), the graded Abhyankar criterion13.8.1--5
(LF19772--19820), the full calibration formula/proof14.5.2--3
(LF20953--21008), and the entire atomic calculation8.9.1--4
(LF10595--10739). The LF windows actually retained may include
adjacent context; essential statements and whole immediate proofs
were read in bounded outputs. The full relative atom formula18.6.2
(LF27607--27671) was read as corroboration, not a premise of the
finite-flat/calibration proof.
No algebraically closed standing hypothesis from14.6 is imported.
The real-plane coefficient one was computed above rather than
inferred from an OCR inequality.

The continuous products and bounded Green foundations retain the
exact earlier 8.5--8.8 source applications and all their hypotheses;
the complete fresh point construction and energy argument are
given here. GR is the full17September2025 v2 source, bound by
\texttt{3a221b86116604c950db07e4c60ab184701f0624deebabd92502e5c81871436b}.
Lemma6.3 and its whole immediate proof were read at LF1673--1780.
Its deeper Temkin/Ducros reduction proofs, basic Berkovich
relative-interior/dimension and completed-residue facts, graded
field transcendence, and the analytic quotient construction are
disclosed conventional foundations. There is no claim of an
entire new recursive source audit. All full primary bodies,
literal units, pixels and book contexts remain LOCALONLY.

Original P16.7 still requires actual separated energy-sublevel
compactness/representability, envelopes or a proved variational
replacement, arbitrary-support equation existence at the stated
full bounded-dual criterion, and full generated pointwise uniqueness.
This paper proves the previously missing actual point-finiteness
and coordinate lower bounds. It neither closes the original nor
changes any accepted count. The protected P16.4 editor remains open
and unchanged; new work is not sent as hot feedback to active
native campaigns. All claims remain UNREVIEWED.

\end{document}
